\chapter{绪论}

绪论部分主要围绕课题提出背景、研究价值、相关研究现状以及本文拟解决的问题展开论述，其作用在于说明“为什么要做工程图纸相似检索”“现有方法还缺什么”以及“本文准备如何组织研究工作”。通过本章的铺垫，可以为后续系统设计、方法构建与实验分析提供清晰的问题定义和研究主线。

\section{研究背景}

随着机械制造、工艺设计和产品研发过程逐渐数字化，企业和科研机构积累了大量历史工程图纸。工程图纸不仅记录零件结构与装配关系，还包含图号、材料、比例、技术要求等重要信息。对于工程人员而言，从历史图库中快速检索出与当前图纸相似的图纸，能够显著提高设计复用效率，减少重复设计和试错成本。

传统工程图纸管理方式多依赖人工目录分类、图号检索或关键字搜索。这些方式虽然适合“已知图号”的场景，但在“只知道当前图纸外观、希望找到结构相近历史图纸”的任务中效果有限。尤其是在大量零件图、装配图和工序图混合存储的情况下，仅依赖文本信息难以完成高质量检索。

近年来，深度学习特别是视觉语义嵌入模型的发展为图像检索提供了新的思路。CLIP 作为典型的视觉--语言预训练模型，具备较强的跨任务泛化能力~\citep{radford2021clip}。然而工程图纸不同于自然图像，其主要信息由线条布局、几何轮廓、标题栏和尺寸标注构成，颜色和纹理特征极弱。因此，将通用视觉模型直接用于工程图纸检索往往会忽略结构主导信息。围绕这一问题开展有针对性的模型改进，具有明确的学术价值与工程意义。

\section{研究意义}

本文研究的意义主要体现在以下几个方面。

\begin{enumerate}
  \item 在应用层面，工程图纸相似检索有助于提升企业历史图纸复用效率和设计标准化水平。
  \item 在方法层面，工程图纸场景为 CLIP 等通用视觉模型的工业迁移提供了具有代表性的研究对象。
  \item 在工程层面，构建完整的检索服务、实验评测与可解释性工具链，有助于将算法研究转化为可落地系统。
  \item 在论文层面，围绕双主线检索框架、结构重排、标题栏字段级 OCR 辅助和 Grad-CAM 分析形成了较完整的方法设计与实验支撑。
\end{enumerate}

\section{研究目标与主要指标}

结合任务书与开题报告中对课题目标的要求，本文拟完成一个面向机械设计场景的工程图纸检索系统，并在算法效果与系统性能两个层面达到可验证的毕业设计目标。具体而言，本文希望在完成数据集清理与系统实现的基础上，重点达成以下目标：

\begin{enumerate}
  \item 构建面向工程图纸场景的专用图库，形成覆盖多类别机械零部件与装配图的标准化数据基础；
  \item 建立基于预训练模型迁移的视觉检索框架，并围绕工程图纸场景完成主体结构增强、结构重排与文本辅助检索的改进设计；
  \item 通过正式实验验证不同技术模块的作用差异，形成可复查、可比较的实验结论；
  \item 实现可运行的原型系统，使图纸上传、特征提取、相似检索、结果展示和图库管理形成完整闭环。
\end{enumerate}

从量化指标角度看，任务书与开题报告对系统提出了较明确的目标要求，即在大规模图纸数据上达到较高的 Top-1 检索准确率与较好的前列排序质量，同时保证端到端响应时间满足工程应用需求。本文在实际研究过程中延续了这一目标导向，并进一步采用 Recall@K、mAP@10、nDCG@10、FT、ST 和 ANMRR 等多项指标对系统进行综合评价，从而使实验结果不仅能够体现单点命中能力，也能够体现候选排序整体质量。

\section{国内外研究现状}

图像检索研究大体经历了手工特征方法、卷积神经网络方法以及视觉 Transformer 和多模态预训练模型方法的发展阶段。传统手工特征方法依赖边缘、纹理和形状描述子，但对复杂工程图纸的鲁棒性不足。深度特征方法通过学习统一嵌入空间提高了检索能力，但自然图像预训练特征在工程图纸场景中容易受模板边框和局部噪声影响~\citep{zhang2023cbirsurvey,liu2024matchingreview}。

近年来，CLIP 等多模态预训练模型在图像理解和检索任务中表现优异~\citep{radford2021clip,han2024crossmodalreview}，相关工作也证明了其在草图检索等结构主导任务中的潜力~\citep{sain2023clipforallthings}。与此同时，参数高效微调技术，如 CLIP-Adapter~\citep{gao2024clipadapter} 与 CoOp~\citep{zhou2022coop}，为将大模型迁移到垂直场景提供了可行路径。

在 CAD 图形与工程图纸检索方向，已有研究尝试结合轮廓、结构关系、图号和文本信息进行混合检索~\citep{mahajan2025orthocad,heidari2025geometriccad}。然而，多数工作要么依赖复杂规则，要么缺乏完整实验评测与系统实现验证。特别是如何将标题栏字段作为强语义信息融入视觉检索结果，并在不破坏召回性能的前提下改善前排排序，仍值得深入研究。

\section{研究内容与创新点}

本文围绕工程图纸相似检索任务，主要开展以下工作：

\begin{enumerate}
  \item 构建基于 CLIP 视觉嵌入的工程图纸检索基线。
  \item 将方法路线重构为 YOLO 清洗路线和无清洗注意力式视觉增强路线，并围绕两条主线组织正式对比实验。
  \item 设计掩码引导的局部主体增强机制与结构描述子重排序方法，提高结构相似检索能力。
  \item 设计 OCR 文本辅助方案，并进一步实现标题栏字段级多模态排序增强。
  \item 实现完整的系统原型、正式实验脚本、误检分析和 Grad-CAM 可解释性工具。
\end{enumerate}

本文的主要创新点如下：

\begin{enumerate}
  \item 提出了一种面向工程图纸的双主线检索组织方式，将原有串行堆叠流程重构为 YOLO 清洗路线与无清洗注意力式视觉增强路线，从而能够更清楚地比较显式清洗与编码阶段主体增强两种思路。
  \item 提出了一种掩码引导的局部主体增强与轻量结构描述子重排序组合方法，将主体区域强化和结构先验用于向量召回后的候选精排。
  \item 提出了一种标题栏字段级 OCR 多模态辅助方法，从标题栏中抽取图号、零件名、材料和比例等字段，作为候选重排序修正信号；同时通过正式实验如实验证其当前增益并不稳定。
  \item 提出并实现了基于 Grad-CAM 的成对检索可视化工具，用于解释模型在判定两张图纸相似时的关注区域。
\end{enumerate}

\section{论文结构安排}

全文共分为七章。第 1 章介绍研究背景、研究意义与创新点；第 2 章介绍相关理论与关键技术；第 3 章给出系统总体设计；第 4 章重点介绍改进检索方法；第 5 章介绍实验设计与结果分析；第 6 章介绍系统实现；第 7 章对全文进行总结并展望未来工作。

综上，本章完成了课题背景、研究问题与论文整体安排的说明，明确了本文以工程图纸相似检索为核心任务，以双主线方法设计和系统实现验证为主要研究路径。后续章节将在此基础上进一步展开相关理论、系统结构、方法细节与实验结果分析。
